\begin{center}
\LARGE
``NORT'' Package for Treating Multivariate Power Series
                        and Applications

\end{center}

\begin{center}
\Large
Victor Edneral

\end{center}

\noindent
Date:  July 19th (Friday)\\
Time:  14:25-14:45\\

\begin{center}
\large
Abstract
\end{center}

In the report discussed applications of the NORT package for problems of creation multivariate
power series which approximate periodic solutions of nonlinear autonomous ODEs
(Edneral,1993), trajectories of periodic orbits of such systems, Lyapunov's values for small limit
circles of a cubic plane system of general type and for solving of systems of equation over
truncated formal power series (Edneral,1995), (Vasiliev,Edneral,1995). 

Originally the NORT package was created for building Poincare -- Dulac normal forms
(Edneral,Khrustalev,1985,1992) till fixed order. It is written in Standard LISP and contains
procedures for treating truncated multivariate power series of an arbitrary dimension. Besides of
procedures for arithmetic operations with such series there are special procedures for the creation
of normal forms and procedures for substitutions, calculations of roots, for differentiating, for
printing and for inverting multivariate power series. The package can be used as a separate unit
under a Standard Lisp environment or as a REDUCE package. 

Complex numerical coefficients of the truncated power series may be evaluated in three different
arithmetics: precision rational, floating point and approximate rational (Rodionov A.Ya,in
progress), in according to the user's choice. Separately user may choose a representation of
numbers (in a floating point or in a rational form) for output. This output has a REDUCE
acceptable form, but REDUCE's internal procedures are not used by the NORT. 

The program uses a recurrent polynomial representation separated in each order of the stored
series. There is a compact enough use of memory. So for 8 - dimensional problem we managed to
treat pieces of series with 150000 elements on computer with 64 Mbyte of RAM. 

\bigskip
\noindent 
{\bf References}
\begin{enumerate}

\item
Edneral, V.F., Khrustalev, O.A. (1985) The Normalizing Transformation for Nonlinear Systems of
ODE. The Realization of the Algorithm; Proceedings of the International Conference on Computer
Algebra and its Application in Theoretical Physics (USSR, Dubna, September 1985). Dubna: JINR
publ., 219--224. In Russian. 


\item
Edneral, V.F., Khrustalev, O.A. (1992) Program for Recasting ODE Systems in the Normal Form;
``Programmirovanie'' (Sov. J. of Programming), \# 5, 73--80. In Russian. 


\item
Edneral, V.F. (1993) Computer Generation of Normalizing Transformation for Systems of
Nonlinear ODE; Proceedings of the 1993 International Symposium on Symbolic and Algebraic
Computation (Kiev, Ukraine, July 6-8, 1993). N.Y., ACM Press, edited by M.Bronstein, pp. 14--
19. 


\item
Vasiliev N.N., Edneral V.F., (1995) About Constructing Canonical Basis for the Ring of Formal
Power Series; The Abstracts of the conference ``Computer Methods of Celestial Mechanics
(S.-Petersburg, Russia, October 17--20, 1995)'', pp. 72--73. In Russian. 


\item
Edneral V.F., (1995) To the Question of Multidimensional Power Series Solving; The Abstracts of
the PoSSo-95 ``PoSSo Open Workshop on Applications of PoSSo and Real Solving, (Iraklio,
Greece, June 7--10, 1995)'', p. 34. 


\end{enumerate}